\section*{Bab \ref{ch:b3:mantle-convection} --- Konveksi Mantel dan Gaya Penggerak Lempeng}

\begin{solution}{exo:b3:mantle-convection:1}
$\mathrm{Ra} = 3500\times9.9\times\num{3e-5}\times1000\times(\num{6.7e5})^3/
(10^{-6}\times\num{5e20}) = \num{6.25e5}$, sekitar 950 kali \omterm{def:b3:mantle-convection:critical}{bilangan Rayleigh
kritis} 657.5: lapisan berkonveksi kuat.
\end{solution}

\begin{solution}{exo:b3:mantle-convection:2}
$\tau_M = 10^{21}/\num{7.7e10} = \qty{1.30e10}{s}$, sekitar \qty{412}{yr}. Pasang
surut, yang berlangsung berjam-jam, bertemu mantel elastis; pencairan lembaran es,
selama ribuan tahun, dan tenggelamnya lempengan selama jutaan tahun, bertemu mantel
viskos.
\end{solution}

\begin{solution}{exo:b3:mantle-convection:3}
$\tau = 4\pi\times10^{21}/(3300\times9.8\lambda)$: \qty{3.89e11}{s} =
\qty{12.3}{ka} untuk \qty{1000}{km}, dan \qty{4.1}{ka} untuk \qty{3000}{km}. Beban
yang lebih luas pulih lebih cepat.
\end{solution}

\begin{solution}{exo:b3:mantle-convection:4}
$g\rho_m\alpha(T_m - T_0)\kappa t = 9.8\times3300\times\num{3.1e-5}\times1450
\times\num{0.8e-6}\times50\times\num{3.156e13} = \qty{1.835e12}{N.m^{-1}}$; faktor
dalam kurung bernilai $1 + 2\times3300\times\num{3.1e-5}\times1450/(\pi\times2270) = 1.0416$,
sehingga $F_{rp} = \qty{1.91e12}{N.m^{-1}}$, 24 kali lebih kecil daripada tarikan lempengan.
\end{solution}

\begin{solution}{exo:b3:mantle-convection:5}
(a) $v = 3300\times9.8\times\num{3e-5}\times300\times(\num{e4})^2/10^{21} =
\qty{2.91e-11}{m.s^{-1}}$; waktu naik sepanjang \qty{10}{km} adalah $d/v =
\qty{3.44e14}{s}$, \qty{10.9}{Ma}, waktu difusi $d^2/\kappa =
\qty{e14}{s}$, \qty{3.2}{Ma}. Rasio keduanya adalah $(d^2/\kappa)/(d/v) =
vd/\kappa = \rho g\alpha\Delta Td^3/(\kappa\eta) = \mathrm{Ra} = 0.29$. (b) Ia
mendingin lebih dulu: gumpalan kecil tidak berkonveksi. (c) $\mathrm{Ra} \propto d^3$, sehingga $d =
10\times(657.5/0.291)^{1/3} = \qty{131}{km}$.
\end{solution}

\begin{solution}{exo:b3:mantle-convection:6}
$2\mathrm{Ra}_c = 1315$ dan $\mathrm{Ra}/2\mathrm{Ra}_c = 7604$, yang akar pangkat tiganya
adalah 19.66. $\delta = 2891/19.66 = \qty{147}{km}$, $\mathrm{Nu} = 19.66/2 =
9.8$, $u = (10^{-6}/\num{2.891e6})\times19.66^2 = \qty{1.34e-10}{m.s^{-1}} =
\qty{0.42}{cm/yr}$, dan $q = 9.8\times3.1\times1450/\num{2.891e6} =
\qty{15}{mW.m^{-2}}$.
\end{solution}

\begin{solution}{exo:b3:mantle-convection:7}
$u = 0.09/\num{3.156e7} = \qty{2.85e-9}{m.s^{-1}}$ dan $B = 2270\times\num{8e8}
\times\num{2.85e-9} = \qty{5.18e3}{kg.s^{-1}} = \qty{5.2}{Mg/s}$, sekitar
tiga perlima milik Hawaii. Fluks panas adalah $1170\times5178/\num{3.1e-5} =
\qty{1.95e11}{W}$, \qty{0.2}{TW}, atau \qty{0.42}{\percent} dari kehilangan
panas Bumi.
\end{solution}

\begin{solution}{exo:b3:mantle-convection:8}
$v = 2\times30\times9.8\times(\num{4e5})^2/(9\times10^{21}) =
\qty{1.05e-8}{m.s^{-1}} = \qty{0.33}{m/yr}$; \qty{2000}{km} butuh
$\num{2e6}/\num{1.05e-8} = \qty{1.9e14}{s}$, sekitar \qty{6}{Ma}.
\end{solution}

\begin{solution}{exo:b3:mantle-convection:9}
$w = 300e^{-10/4.4} = \qty{30.9}{m}$ dan $\dot w = w/\tau = 30.9/4400 =
\qty{7.0}{mm/yr}$, dua pertiga laju yang diukur dekat Umeå: setengah-ruang seragam
dengan $\tau$ ini berada pada rentang yang tepat, tetapi pantulan sebenarnya
lebih cepat atau cekungannya lebih dalam.
\end{solution}

\begin{solution}{exo:b3:mantle-convection:10}
(a) Dengan $w = 2\rho_m\alpha\Delta T\sqrt{\kappa t/\pi}/(\rho_m - \rho_w)$,
\[
  \tfrac12(\rho_m - \rho_w)gw^2 = \frac{2g\rho_m^2\alpha^2\Delta T^2\kappa t}
  {\pi(\rho_m - \rho_w)} = g\rho_m\alpha\Delta T\kappa t\times
  \frac{2\rho_m\alpha\Delta T}{\pi(\rho_m - \rho_w)} .
\] (b) Dengan integrasi parsial,
$\int_0^\infty x\operatorname{erfc}x\,\mathrm dx = [\tfrac{x^2}2\operatorname{erfc}
x]_0^\infty + \int_0^\infty\tfrac{x^2}2\,\tfrac2{\sqrt\pi}e^{-x^2}\,\mathrm dx =
0 + \tfrac1{\sqrt\pi}\cdot\tfrac{\sqrt\pi}4 = \tfrac14$. (c) Pada \qty{100}{Ma},
$\rho_mg\alpha\int z\theta = \qty{3.67e12}{N.m^{-1}}$ dan, dengan $w =
\qty{3705}{m}$, $\tfrac12(\rho_m - \rho_w)gw^2 = \qty{1.53e11}{N.m^{-1}}$: suku
air hanya \qty{4}{\percent} dari gaya total.
\end{solution}

\begin{solution}{exo:b3:mantle-convection:11}
$\ln f = 3\ln(\pi^2 + k^2) - 2\ln k$, yang turunannya $6k/(\pi^2 + k^2) - 2/k$
bernilai nol ketika $3k^2 = \pi^2 + k^2$, $k^2 = \pi^2/2$. Di sana $f = (3\pi^2/2)^3/
(\pi^2/2) = 27\pi^4/4$. Panjang gelombang $2\pi/k = 2\sqrt2\,d = 2.83\,d$ memuat
sepasang sel yang berputar berlawanan arah, masing-masing selebar $\sqrt2\,d = 1.41\,d$.
\end{solution}

\begin{solution}{exo:b3:mantle-convection:12}
(a) $u = \qty{1.90e-9}{m.s^{-1}}$, $\eta u/h = 10^{20}\times\num{1.90e-9}/
\num{1.5e5} = \qty{1.27}{MPa}$, dan sepanjang \qty{6000}{km}, \qty{7.6e12}{N.m^{-1}}.
(b) Dua kali dorongan punggung dan seperenam tarikan lempengan: seretan dapat menahan
dorongan punggung saja, tetapi tidak tarikan lempengan. (c) $\eta uL/h = \num{3.8e12} +
\num{4.6e12}$ memberikan $u = \num{8.4e12}\times\num{1.5e5}/(\num{6e6}\times10^{20})
= \qty{2.1e-9}{m.s^{-1}} = \qty{6.6}{cm/yr}$, laju lempeng samudra
yang cepat.
\end{solution}

\begin{solution}{pb:b3:mantle-convection:1}
\textbf{1.} $p = 1000\times9.8\times300 = \qty{2.94}{MPa}$.

\textbf{2.} $w = \rho_wh/\rho_m = 300\times1000/3300 = \qty{90.9}{m}$. Kerak
bergerak turun sebagai satu blok; yang mengalir keluar dari bawah beban, lalu kembali,
adalah mantel, sehingga kolom diseimbangkan oleh mantel yang terdesak di kedalaman.

\textbf{3.} $\pi\times(\num{1.25e5})^2 = \qty{4.91e10}{m^2}$, dan air
menimbang $\num{4.91e10}\times300\times1000 = \qty{1.47e16}{kg}$.

\textbf{4.} Garis pantai terbentuk datar, ketika dasar paling turun di
pusat; sejak itu pusat naik lebih banyak daripada tepi, sehingga garis kini
naik menuju pusat, sebesar \qty{77.2}{m} sepanjang \qty{125}{km}:
\qty{0.62}{m/km}.

\textbf{5.} $\lambda/2 = \qty{250}{km}$, sehingga $\lambda = \qty{500}{km}$ dan $k =
2\pi/\num{5e5} = \qty{1.257e-5}{m^{-1}}$.

\textbf{6.} $1/k = \lambda/2\pi = \qty{79.6}{km}$.

\textbf{7.} $w = 90.9 - \text{pengangkatan}$: 79.4, 69.3, 52.9, 40.4, 26.9 dan
\qty{13.7}{m}.

\textbf{8.} $w = w_0e^{-t/\tau}$ memberikan $\ln w = \ln w_0 - t/\tau$, sebuah garis lurus
dengan kemiringan $-1/\tau$.

\textbf{9.} $\tau = (14 - 1)/\ln(79.4/13.7) = 13/1.7571 = \qty{7.40}{ka}$.

\textbf{10.} $w(6) = 79.4\,e^{-5/7.40} = \qty{40.4}{m}$, pengangkatan sebesar
\qty{50.5}{m}, sesuai dengan tabel.

\textbf{11.} \qty{13.7}{m} tersisa; $\dot w = 13.7/7400 = \qty{1.85}{mm/yr}$.

\textbf{12.} Seluruh \qty{13.7}{m}, didekati makin lama makin lambat.

\textbf{13.} Laju turun mengikuti $e^{-t/\tau}$: $t = 7.40\ln(1.85/0.5) =
7.40\times1.309 = \qty{9.7}{ka}$, sekitar 9700 tahun lagi.

\textbf{14.} $\tau = 7399\times\num{3.156e7} = \qty{2.335e11}{s}$, dan $\eta =
\tau\rho_mg\lambda/(4\pi) = \num{2.335e11}\times3300\times9.8\times\num{5e5}/
12.566 = \qty{3.0e20}{Pa.s}$.

\textbf{15.} Tepat di bawah ujung bawah rentang Fennoscandia. Aliran
di bawah beban dengan panjang gelombang \qty{500}{km} meluruh dalam \qty{80}{km}: 
pantulan mengambil sampel mantel terdangkal dan terlemah, di bawah dan sebagian di dalam
yang di Bumi nyata adalah litosfer.

\textbf{16.} $\tau = 4\pi\times\num{3.0e20}/(3300\times9.8\times\num{3e6}) =
\qty{3.9e10}{s} = \qty{1.2}{ka}$. Di bawah beban selebar itu aliran mencapai
hingga \qty{477}{km} lebih dalam, menuju mantel bawah yang sepuluh kali lebih viskos,
yang memperlambatnya.

\textbf{17.} Litosfer menahan sebagian beban sempit melalui lenturan: depresi pada
kesetimbangan lebih dangkal dan lebih lebar daripada \qty{90.9}{m} isostatik, dan sebagian
pantulan itu bersifat elastik dan seketika. Setengah-ruang kental seragam yang dicocokkan dengan
pengangkatan kemudian memberikan viskositas yang mencampur kedua efek.

\textbf{18.} $\tau_M = \num{3.0e20}/\num{7.7e10} = \qty{3.9e9}{s}$, sekitar
\qty{124}{yr}, enam puluh kali lebih pendek daripada $\tau$: selama ribuan tahun
mantel bereaksi seperti fluida kental.

\textbf{19.} $1 - e^{-50/7.40} = 0.9988$: \qty{99.9}{\percent} dari
depresi kesetimbangan telah tercapai, dan asumsi ini berlaku.

\textbf{20.} $\eta \propto \lambda$: memangkas $\lambda$ menjadi separuh memangkas $\eta$ menjadi separuh, dan
$\log_{10}\eta$ turun sebesar $\log_{10}2 = 0.30$.

\textbf{21.} $\mathrm{Ra} = 3700\times9.9\times\num{3.1e-5}\times1450\times
(\num{6.7e5})^3/(10^{-6}\times\num{3.0e20}) = \num{1.65e6}$.

\textbf{22.} $\num{1.65e6}/657.5 \approx 2500$ kali kritis.

\textbf{23.} $(\mathrm{Ra}/2\mathrm{Ra}_c)^{1/3} = (1254)^{1/3} = 10.78$: $\delta
= 670/10.78 = \qty{62}{km}$ dan $\mathrm{Nu} = 5.4$.

\textbf{24.} $u = (10^{-6}/\num{6.7e5})\times10.78^2 = \qty{1.73e-10}{m.s^{-1}}
= \qty{0.55}{cm/yr}$, dua sampai empat belas kali lebih lambat daripada lempeng.

\textbf{25.} $q = 5.4\times3.1\times1450/\num{6.7e5} = \qty{36}{mW.m^{-2}}$. Lapisan
ini mengabaikan panas yang dihasilkan dalam mantel dan panas dari
inti serta mantel bawah; dan penskalaan menghilangkan konstanta berorde satu,
sedangkan lempeng, yang mendingin sebagai setengah-ruang dari punggungan, kehilangan jauh lebih banyak panas pada
dasar samudra muda daripada lapisan batas rata-rata.

\textbf{26.} $\eta = \qty{3.0e20}{Pa.s}$: $\log_{10}\eta =
\textbf{\num[text-series-to-math=true]{20.48}}$.
\end{solution}
